\chapter{The finite tile-count model}

Write \([n]=\{1,\ldots,n\}\).  All coordinates in the mathematical
description are one-based, even though an implementation may store them with
zero-based finite indices.

\begin{definition}\leanok
[Count vectors]
  \label{def:counts}
  \lean{Ninegates.Counts}
  % SOURCE: half-formal.qmd, “本文范围与记号约定” (read from references/half-formal.qmd)
  % SOURCE QUOTE: “一个清一色牌型用重数向量 $c=(c_1,\ldots,c_n)$ 表示，$c_i$ 是点数 $i$ 的张数。”
  \textit{Source: half-formal.qmd, notation conventions.}
  A count vector of length \(n\) is a function \(c:[n]\to\mathbb N\).
  Its value \(c_i\) is the multiplicity of rank \(i\).
\end{definition}

\begin{definition}\leanok
[Basis count vector]
  \label{def:basis}
  \lean{Ninegates.basis}
  \uses{def:counts}
  % SOURCE: half-formal.qmd, “本文范围与记号约定” (read from references/half-formal.qmd)
  % SOURCE QUOTE: “$e_i$ 表示第 $i$ 个标准基向量，即第 $i$ 个坐标为 $1$，其余坐标为 $0$。”
  \textit{Source: half-formal.qmd, notation conventions.}
  For \(i\in[n]\), the vector \(e_i\) has value \(1\) at \(i\) and
  value \(0\) at every other rank.
\end{definition}

\begin{definition}\leanok
[Total tile count]
  \label{def:total}
  \lean{Ninegates.total}
  \uses{def:counts}
  The total number of tiles in \(c\) is
  \[
    |c|=\sum_{i\in[n]}c_i.
  \]
\end{definition}

\begin{definition}\leanok
[Removing a pair]
  \label{def:remove_pair}
  \lean{Ninegates.removePair}
  \uses{def:counts,def:basis}
  If \(c_p\geq 2\), removing a pair at \(p\) means coordinatewise
  natural-number subtraction \(c-2e_p\).  The definition is available for
  every \(c,p\); the inequality is imposed by statements that use it.
\end{definition}

\begin{definition}\leanok
[Meldability]
  \label{def:meldable}
  \lean{Ninegates.Meldable}
  \uses{def:counts}
  % SOURCE: half-formal.qmd, #def-meldable (read from references/half-formal.qmd)
  % SOURCE QUOTE: “设 $d=(d_1,\ldots,d_m)$ 是非负整数向量。称 $d$ **可面子分解**，如果存在非负整数 $t_1,\ldots,t_m$ 和 $r_1,\ldots,r_{m-2}$，使得对每个 $i\in\{1,\ldots,m\}$ 都有 $$ d_i=3t_i+r_i+r_{i-1}+r_{i-2}. \tag{1} $$ 这里约定，凡是下标不属于 $\{1,\ldots,m-2\}$ 的 $r_j$ 均为 $0$。”
  \textit{Source: half-formal.qmd, Definition \texttt{\#def-meldable}.}
  A count vector \(d\) of length \(m\) is meldable if there are
  nonnegative triplet counts \(t_i\) and run-start counts \(r_i\) such that
  \[
    d_i=3t_i+r_i+r_{i-1}+r_{i-2}
  \]
  for every \(i\in[m]\).  Run counts outside \(1,\ldots,m-2\) are
  interpreted as zero.  Thus no run may extend beyond the right endpoint.
\end{definition}

\begin{definition}\leanok
[Complete hand]
  \label{def:complete}
  \lean{Ninegates.Complete}
  \uses{def:counts,def:remove_pair,def:meldable}
  % SOURCE: half-formal.qmd, #def-complete (read from references/half-formal.qmd)
  % SOURCE QUOTE: “设 $w=(w_1,\ldots,w_n)$ 是非负整数向量。称 $w$ 是**完整牌型**，如果存在雀头点数 $p\in\{1,\ldots,n\}$，满足 $w_p\ge2$，并且 $w-2e_p$ 可面子分解。”
  \textit{Source: half-formal.qmd, Definition \texttt{\#def-complete}.}
  A count vector \(w\) is complete if some \(p\in[n]\) satisfies
  \(w_p\geq2\) and \(w-2e_p\) is meldable.
\end{definition}

\begin{definition}\leanok
[Live all-rank wait]
  \label{def:live_wait}
  \lean{Ninegates.LiveWait}
  \uses{def:counts,def:basis,def:total,def:complete}
  % SOURCE: half-formal.qmd, #def-live-wait (read from references/half-formal.qmd)
  % SOURCE QUOTE: “一个向量 $c=(c_1,\ldots,c_n)$ 称为广义 $n$ 阶麻将中的**非空听 $n$ 面听清一色牌型**，如果满足： 1. 对每个 $i\in\{1,\ldots,n\}$，$$0\le c_i\le3;$$ 2. 摸牌前的总张数为 $$\sum_{i=1}^n c_i=n+4;$$ 3. 对每个进张 $x\in\{1,\ldots,n\}$，向量 $$c+e_x$$ 都是完整牌型。”
  \textit{Source: half-formal.qmd, Definition \texttt{\#def-live-wait}.}
  A count vector \(c\) is a live all-rank wait if:
  every coordinate is at most \(3\), its total is \(n+4\), and
  \(c+e_x\) is complete for every \(x\in[n]\).
\end{definition}

\begin{definition}\leanok
[The standard hand]
  \label{def:standard}
  \lean{Ninegates.standard}
  \uses{def:counts}
  % SOURCE: half-formal.qmd, “把任意反例压缩到有限规模” (read from references/half-formal.qmd)
  % SOURCE QUOTE: “称一个非空听的 $n$ 面听牌型为**非标准牌型**，如果它不等于 $G_n=(3,1,\ldots,1,3)$。”
  \textit{Source: half-formal.qmd, finite-size reduction section.}
  The standard count vector \(G_n\) has multiplicity \(3\) at the two
  endpoints and multiplicity \(1\) at every internal rank.  The definition
  is made uniformly for all \(n\); the structural theorems only use
  \(n\geq6\).
\end{definition}
